\documentclass[11pt,a4paper]{article}

\usepackage[T1]{fontenc}
\usepackage[utf8]{inputenc}
\usepackage[ngerman]{babel}

\usepackage{amsmath,amssymb,amsthm}
\usepackage{graphicx}
\usepackage{hyperref}
\usepackage{geometry}

\geometry{margin=2.5cm}

\title{Geschlossene Teleportationszyklen, Horizontprojektionen und scheinbarer Energietransport\\
	Ein konzeptioneller Rahmen}

\author{Thomas Hoffbauer}

\date{\today}

\begin{document}
	
	\maketitle
	
	\begin{abstract}
		
		Wir entwickeln einen konzeptionellen Rahmen, der das No-Cloning-Theorem der Quantenmechanik, zyklische Holonomiestrukturen und beobachterabhängige Horizontprojektionen miteinander verbindet.
		
		Die zentrale Idee besteht darin, physikalischen Transport nicht als Bewegung persistenter Objekte durch den Raum zu interpretieren, sondern als Folge von Zustandsrekonstruktionen, analog zur Quantenteleportation.
		
		Es wird gezeigt, dass geschlossene Zustandszyklen global energieerhaltend sein können, während für beschleunigte Beobachter aufgrund von Rindler-Horizonten ein effektiver lokaler Energiefluss erscheint.
		
		Das Modell liefert eine neue Perspektive auf die Beziehung zwischen Information, Holonomie, Horizonten und den klassischen Fragen der Einweg- und Zweiweg-Lichtgeschwindigkeit.
		
	\end{abstract}
	
	\section{Einleitung}
	
	Die Quantenteleportation zeigt, dass unbekannte Quantenzustände nicht kopiert werden können.
	
	Stattdessen verschwindet der Zustand an einem Ort und wird an einem anderen Ort rekonstruiert.
	
	Formal:
	
	\[
	|\psi\rangle_A |0\rangle_B
	\longrightarrow
	|0\rangle_A |\psi\rangle_B.
	\]
	
	Dies legt eine alternative Interpretation nahe:
	
	\begin{quote}
		Transport entspricht nicht notwendigerweise der Bewegung eines Objektes, sondern kann als Rekonstruktion eines Zustandes entlang einer Kette verschränkter Freiheitsgrade verstanden werden.
	\end{quote}
	
	In dieser Arbeit untersuchen wir die Konsequenzen geschlossener Zustandszyklen unter dieser Annahme.
	
	\section{Das No-Cloning-Transportprinzip}
	
	Sei
	
	\[
	\mathcal H
	\]
	
	der Zustandsraum.
	
	Nach dem No-Cloning-Theorem existiert kein universeller Operator
	
	\[
	|\psi\rangle |0\rangle
	\rightarrow
	|\psi\rangle |\psi\rangle.
	\]
	
	Erlaubt ist dagegen eine Zustandsübertragung:
	
	\[
	T_{A\to B}
	\left(
	|\psi\rangle_A |0\rangle_B
	\right)
	=
	|0\rangle_A |\psi\rangle_B.
	\]
	
	Es entsteht kein Duplikat.
	
	Der ursprüngliche Zustand verschwindet.
	
	\section{Der EABC-Zyklus}
	
	Wir betrachten vier Zustandsklassen
	
	\[
	E,\;A,\;B,\;C.
	\]
	
	Ein Zyklusoperator sei definiert durch
	
	\[
	R:
	E\rightarrow A
	\rightarrow B
	\rightarrow C
	\rightarrow E.
	\]
	
	Es gilt
	
	\[
	R^4=I.
	\]
	
	Der vollständige Zyklus besteht aus den vier Teleportationsschritten
	
	\[
	E\rightarrow A,
	\quad
	A\rightarrow B,
	\quad
	B\rightarrow C,
	\quad
	C\rightarrow E.
	\]
	
	Der Gesamtoperator lautet
	
	\[
	\mathcal T
	=
	T_{C\to E}
	T_{B\to C}
	T_{A\to B}
	T_{E\to A}.
	\]
	
	Nach einem Umlauf gilt
	
	\[
	\mathcal T|\psi\rangle
	=
	e^{i\Gamma}
	|\psi\rangle.
	\]
	
	Die Größe
	
	\[
	\Gamma
	\]
	
	bezeichnet die Holonomie des Zyklus.
	
	\section{Holonomie geschlossener Zustandszyklen}
	
	Für einen geschlossenen Zyklus gilt
	
	\[
	\mathcal T^\dagger \mathcal T=I.
	\]
	
	Der Zustand kehrt auf sich selbst zurück.
	
	Allerdings kann eine nichtverschwindende geometrische Phase auftreten:
	
	\[
	\Gamma \neq 0.
	\]
	
	Dies erinnert an
	
	\begin{itemize}
		\item Berry-Phasen,
		\item Aharonov-Bohm-Effekte,
		\item Holonomien in Eichtheorien.
	\end{itemize}
	
	Der Zyklus besitzt damit eine globale Erinnerung an seinen Weg durch den Zustandsraum.
	
	\section{Rindler-Horizonte als Projektionen}
	
	Ein beschleunigter Beobachter besitzt nur Zugang zu einem Teil der Raumzeit.
	
	Sei
	
	\[
	P_H
	\]
	
	die Projektion auf den zugänglichen Rindler-Keil.
	
	Dann sieht der Beobachter nicht den vollständigen Zyklus
	
	\[
	R^4,
	\]
	
	sondern lediglich
	
	\[
	P_H R^4 P_H.
	\]
	
	Im Allgemeinen gilt
	
	\[
	P_H R^4 P_H
	\neq
	P_H.
	\]
	
	Wir definieren den Horizontrest
	
	\[
	J_H
	=
	P_H
	-
	P_HR^4P_H.
	\]
	
	Dieser misst die beobachtete Nichtschließung des Zyklus.
	
	\section{Beobachterabhängiger Energietransport}
	
	Sei
	
	\[
	\hat H
	\]
	
	der Hamiltonoperator.
	
	Die für den Beobachter zugängliche Energie lautet
	
	\[
	E_H
	=
	\langle \psi |
	P_H \hat H P_H
	|
	\psi
	\rangle.
	\]
	
	Nach einem Zyklus ergibt sich
	
	\[
	\Delta E_H
	=
	\langle \psi |
	\mathcal T^\dagger
	P_H \hat H P_H
	\mathcal T
	|
	\psi
	\rangle
	-
	\langle \psi |
	P_H \hat H P_H
	|
	\psi
	\rangle.
	\]
	
	Global sei
	
	\[
	[\mathcal T,\hat H]=0.
	\]
	
	Dann bleibt die Gesamtenergie erhalten.
	
	Lokal kann jedoch gelten
	
	\[
	\Delta E_H \neq 0.
	\]
	
	Daraus folgt die zentrale Aussage:
	
	\begin{quote}
		Ein global geschlossener Zustandszyklus kann für einen beschleunigten Beobachter als lokaler Energiefluss erscheinen.
	\end{quote}
	
	\section{Kompensation jenseits des Horizonts}
	
	Sei
	
	\[
	Q_H=I-P_H.
	\]
	
	Dann gilt
	
	\[
	\Delta E_H
	+
	\Delta E_Q
	=
	0.
	\]
	
	Jeder lokale Energiegewinn wird durch einen entsprechenden Verlust im nicht zugänglichen Bereich kompensiert.
	
	Es entsteht somit kein Perpetuum Mobile.
	
	\section{Bezug zur Einweg- und Zweiweg-Lichtgeschwindigkeit}
	
	Ein vollständiger Zyklus entspricht einer Rundreise:
	
	\[
	\oint_\gamma ds.
	\]
	
	Ein Horizont zerlegt den Zyklus in
	
	\[
	\gamma_H
	\]
	
	und
	
	\[
	\gamma_Q.
	\]
	
	Somit:
	
	\[
	\oint_\gamma ds
	=
	\int_{\gamma_H} ds
	+
	\int_{\gamma_Q} ds.
	\]
	
	Dies erinnert strukturell an die Unterscheidung zwischen
	
	\begin{itemize}
		\item Zweiweg-Lichtgeschwindigkeit (messbar),
		\item Einweg-Lichtgeschwindigkeit (synchronisationsabhängig).
	\end{itemize}
	
	Horizontprojektionen können daher als geometrische Verallgemeinerung des Synchronisationsproblems interpretiert werden.
	
	\section{Horizont-Holonomie-Hypothese}
	
	\textbf{Hypothese.}
	
	Sei
	
	\[
	\mathcal T
	\]
	
	ein geschlossener Teleportationszyklus mit
	
	\[
	\mathcal T |\psi\rangle
	=
	e^{i\Gamma}
	|\psi\rangle.
	\]
	
	Sei ferner
	
	\[
	P_H
	\]
	
	eine Horizontprojektion.
	
	Dann kann trotz globaler Energieerhaltung
	
	\[
	\mathcal T^\dagger
	\hat H
	\mathcal T
	=
	\hat H
	\]
	
	lokal gelten
	
	\[
	\mathcal T^\dagger
	P_H\hat H P_H
	\mathcal T
	\neq
	P_H\hat H P_H.
	\]
	
	Dies erzeugt einen beobachterabhängigen effektiven Energiefluss.
	
	\section{Ausblick}
	
	Das vorliegende Papier ist bewusst konzeptionell.
	
	Künftige Arbeiten sollten untersuchen:
	
	\begin{enumerate}
		\item explizite Modelle in der Quantenfeldtheorie auf Rindler-Räumen,
		\item Beziehungen zum Unruh-Effekt,
		\item Verschränkungsentropie an Horizonten,
		\item topologische Pumpen und Holonomietransport,
		\item konkrete Realisierungen von EABC-Zyklen.
	\end{enumerate}
	
	\section*{Schlussbemerkung}
	
	Die hier vorgestellte Horizont-Holonomie-Hypothese verbindet drei fundamentale Konzepte moderner Physik:
	
	\[
	\text{No-Cloning}
	\quad+\quad
	\text{Holonomie}
	\quad+\quad
	\text{Horizontprojektion}.
	\]
	
	Sie legt nahe, dass geschlossene globale Informationszyklen für bestimmte Beobachter als offene Energieflüsse erscheinen können.
	
\end{document}